\documentclass[10pt,a4paper]{amsart}
\usepackage[margin=19mm]{geometry}
\usepackage{amsmath,amssymb,mathtools}
\usepackage[scr=rsfs,cal=euler]{mathalfa}
\usepackage{xfrac}
\usepackage[unicode,hidelinks,bookmarks=false]{hyperref}
\newcommand{\ie}{i.e.\ }
\newcommand{\cf}{cf.\ }
\newcommand{\resp}{respectively\ }
\newcommand{\Subsection}{\subsection}
\providecommand{\nobreakdash}{\nobreak\hbox{-}}
\setlength{\emergencystretch}{2em}
\multlinegap0pt
\newcommand{\R}{\mathbb{R}}
\theoremstyle{plain}
\newtheorem{theorem}{Theorem}[section]
\newtheorem{lemma}[theorem]{Lemma}
\newtheorem{proposition}[theorem]{Proposition}
\newtheorem{corollary}[theorem]{Corollary}
\theoremstyle{definition}
\newtheorem{definition}[theorem]{Definition}
\newtheorem{remark}[theorem]{Remark}
\title[A proof of the Holmes-Holroyd-Ramirez cyclic-product conjecture]{A proof of the Holmes-Holroyd-Ramirez cyclic-product conjecture}
\author{Paata Ivanisvili, Shengtong Zhang}
\date{Source: arXiv:2609.27843v1 (CC BY 4.0)}
\begin{document}
\maketitle

\noindent\emph{Source and licence.} The delivered work is the article shown in the title above, version arXiv:2609.27843v1 (CC BY 4.0), distributed under the Creative Commons Attribution 4.0 International licence, as recorded in the arXiv licence metadata for this exact version and shown on the abstract page saved with this item. The full licence notice, which carries the licence URL and the address of that abstract page, is the bundled file \texttt{licenses/arxiv-2609.27843-CC-BY-4.0.txt}.
\par\smallskip

\noindent\textit{Editorial scope.} Counted unit: the article's Lemma (source label lem:smaller-box) (source0.tex lines 410--418, source label \texttt{lem:smaller-box}): ``Smaller box Let be the steep -trapezoid of sum , and let be a circular Lipschitz- integer sequence with the same sum. Set ' tv(t), M' tv(t). If , then .''. It is delivered together with the article's complete original proof (source0.tex lines 420--494, one proof environment adjacent to the statement, 75 lines), copied line for line from the article's own text. Required context retained uncounted: eq:layer (lines 313--316); lem:clip (lines 319--326); lem:full-amp (lines 337--344); lem:padded (lines 369--386). Editorial formatting only: an AMS \texttt{amsart} preamble that restates the article's own macro definitions and its own theorem declarations, and the theorem environments the retained lines use; the article's own class file is neither shipped nor input; the retained environments are renumbered sequentially in order of appearance, whereas the article prints them with its own numbering; the article's equation numbering is not reproduced and every cross-reference inside the retained text resolves to the wrapper's numbering. No citation occurs inside any retained line, so the wrapper carries no bibliography; All retained bodies are exact copies of the bundled \texttt{sources/211519/source0.tex}, so no omission receipt applies. No asset-inclusion command occurs inside any retained span and the package ships no image asset.


\section*{Retained context: eq:layer (source0.tex lines 313--316)}

\begin{equation}\label{eq:layer}
\sum_{i=1}^kv_i^\downarrow
=k\mu+\sum_{j=1}^\Delta\min(k,n_j(v)).
\end{equation}


\section*{Retained context: lem:clip (source0.tex lines 319--326)}

\begin{lemma}\label{lem:clip}
Let $x,y\in\R^\Delta$ be decreasing and satisfy $x\succ y$. Then, for every $k\ge0$,
\[
\sum_{i=1}^\Delta\min(k,x_i)
\le
\sum_{i=1}^\Delta\min(k,y_i).
\]
\end{lemma}


\section*{Retained context: lem:full-amp (source0.tex lines 337--344)}

\begin{lemma}\label{lem:full-amp}
Let $v$ be a circular Lipschitz-$1$ integer sequence with values in $[\mu,M]$, attaining both endpoints, and satisfying $\sum v=S$. Then
\[
n(v)\succ n(\tau),
\qquad\text{and consequently}\qquad
\tau\succ v.
\]
\end{lemma}


\section*{Retained context: lem:padded (source0.tex lines 369--386)}

\begin{lemma}[Padded comparison]\label{lem:padded}
Let $\alpha,\beta$ be nonnegative integers and $\Delta'=\Delta-\alpha-\beta\ge1$. Write
\[
Q'=Q-m\alpha,
\qquad
H'=\frac{Q'}{\Delta'}-\Delta'+1,
\]
and define
\[
n^\sharp
=\bigl(
\underbrace{m,\dots,m}_{\alpha},
\bigl(H'+2(\Delta'-j)\bigr)_{j=1}^{\Delta'},
\underbrace{0,\dots,0}_{\beta}
\bigr)\in\R^\Delta.
\]
If $H'\ge0$ and $H'+2(\Delta'-1)\le m$, then $n^\sharp\succ n(\tau)$.
\end{lemma}


\section{The counted unit: Lemma (source label lem:smaller-box) and its complete original proof}

\begin{lemma}[Smaller box]\label{lem:smaller-box}
Let $\tau$ be the steep $[\mu,M]$-trapezoid of sum $S$, and let $v$ be a circular Lipschitz-$1$ integer sequence with the same sum. Set
\[
\mu'=\min_tv(t),
\qquad
M'=\max_tv(t).
\]
If $[\mu',M']\subsetneq[\mu,M]$, then $\tau\succ v$.
\end{lemma}

\begin{proof}
Let $\Delta'=M'-\mu'$, $\alpha=\mu'-\mu$, $\beta=M-M'$, and $Q'=S-m\mu'$. If $\Delta'=0$, then $v$ is constant, and every vector with the same sum majorises $v$.

Assume $\Delta'\ge1$ and set
\[
n_j^{\mathrm{loc}}(v)=|\{t:v(t)\ge\mu'+j\}|,
\qquad j=1,\dots,\Delta'.
\]
The two-arc argument in Lemma~\ref{lem:full-amp} gives
\[
n_j^{\mathrm{loc}}(v)-n_{j+1}^{\mathrm{loc}}(v)\ge2,
\qquad j=1,\dots,\Delta'-1.
\]
Since $v$ attains $\mu'$ and $M'$, we have $n_{\Delta'}^{\mathrm{loc}}(v)\ge1$ and $n_1^{\mathrm{loc}}(v)\le m-1$. Therefore, for $j=1,\dots,\Delta'$,
\[
1+2(\Delta'-j)
\le n_j^{\mathrm{loc}}(v)
\le m-1-2(j-1).
\]
Summing these inequalities over $j$ gives
\[
\Delta'^2\le Q'=\sum_{j=1}^{\Delta'}n_j^{\mathrm{loc}}(v)
\le\Delta'(m-\Delta').
\]
Thus, for
\[
H'=\frac{Q'}{\Delta'}-\Delta'+1,
\]
we have $H'\ge1$ and $H'+2(\Delta'-1)\le m-1$.

Let
\[
n^\sharp_{\mathrm{loc}}=(H'+2(\Delta'-1),\dots,H')
\]
and set
\[
d_j=n_j^{\mathrm{loc}}(v)-\bigl(H'+2(\Delta'-j)\bigr),
\qquad j=1,\dots,\Delta'.
\]
Then
\[
d_j-d_{j+1}
=n_j^{\mathrm{loc}}(v)-n_{j+1}^{\mathrm{loc}}(v)-2
\ge0,
\qquad j=1,\dots,\Delta'-1,
\]
so $d_1\ge\cdots\ge d_{\Delta'}$. Moreover,
\[
\sum_{j=1}^{\Delta'}d_j
=Q'-\bigl(\Delta'H'+\Delta'(\Delta'-1)\bigr)
=0
\]
by the definition of $H'$. A decreasing sequence with zero sum has nonnegative partial sums, and hence
\[
n^{\mathrm{loc}}(v)\succ n^\sharp_{\mathrm{loc}}.
\]
Define
\[
\widetilde n(v)
=\bigl(
\underbrace{m,\dots,m}_{\alpha},
 n_1^{\mathrm{loc}}(v),\dots,n_{\Delta'}^{\mathrm{loc}}(v),
\underbrace{0,\dots,0}_{\beta}
\bigr).
\]
Since both local vectors are decreasing and take values in $[0,m]$, adjoining the same initial block of $m$'s and terminal block of $0$'s preserves majorisation. Therefore Lemmas~\ref{lem:clip} and~\ref{lem:padded} give, for every $k$,
\[
\sum_{j=1}^\Delta\min(k,\widetilde n_j(v))
\le
\sum_{j=1}^\Delta\min(k,n_j^\sharp)
\le
\sum_{j=1}^\Delta\min(k,n_j(\tau)).
\]
The vector $\widetilde n(v)$ is the superlevel vector of $v$ measured from the baseline $\mu$. By \eqref{eq:layer}, this is $\tau\succ v$.
\end{proof}

\end{document}
